\documentclass[11pt]{article}

\usepackage[letterpaper,margin=1in]{geometry}
\usepackage{amsmath,amssymb}
\usepackage{booktabs}
\usepackage{longtable}
\usepackage{array}
\usepackage{xcolor}
\usepackage{listings}
\usepackage{url}
\usepackage{hyperref}
\usepackage{titlesec}
\usepackage{parskip}
\usepackage{fancyhdr}
\usepackage{enumitem}

\hypersetup{
  colorlinks=true,
  linkcolor=black,
  urlcolor=blue!50!black,
  citecolor=black,
  pdftitle={KV Cache Engine --- Reference Model API},
  pdfauthor={LonghornSilicon},
  pdfsubject={API Reference, Version 0.1}
}

\titleformat{\section}{\normalfont\Large\bfseries}{\thesection.}{0.6em}{}
\titleformat{\subsection}{\normalfont\large\bfseries}{\thesubsection}{0.6em}{}
\titleformat{\subsubsection}{\normalfont\normalsize\bfseries}{\thesubsubsection}{0.6em}{}

\pagestyle{fancy}
\fancyhf{}
\lhead{\small KV Cache Engine --- Reference Model API}
\rhead{\small v0.1}
\cfoot{\thepage}
\renewcommand{\headrulewidth}{0.4pt}

\definecolor{codebg}{gray}{0.97}
\definecolor{codekey}{rgb}{0.13,0.29,0.53}
\definecolor{codecom}{rgb}{0.25,0.5,0.25}
\lstset{
  basicstyle=\ttfamily\small,
  backgroundcolor=\color{codebg},
  frame=single,
  framerule=0.4pt,
  rulecolor=\color{black!30},
  xleftmargin=0.5em,
  xrightmargin=0.5em,
  aboveskip=0.6em,
  belowskip=0.6em,
  showstringspaces=false,
  breaklines=true,
  keywordstyle=\color{codekey}\bfseries,
  commentstyle=\color{codecom}\itshape,
  numbers=none
}

\setlist{itemsep=1pt,topsep=2pt}

\begin{document}

%==============================================================================
\thispagestyle{empty}
\begin{center}
\vspace*{1.2in}
{\Huge\bfseries KV Cache Engine}\\[0.4em]
{\LARGE\bfseries Reference Model API}\\[1.2em]

{\large \textsc{LonghornSilicon LLM Inference Accelerator}}\\[0.4em]
{\large Block 2 of 4 --- Bit-Accurate C++ / C / Python Models}\\[2em]

\rule{0.7\textwidth}{0.4pt}\\[1em]

\begin{tabular}{rl}
\textbf{Version}:    & 0.1 (first public draft) \\
\textbf{Date}:       & 2026-05-14 \\
\textbf{Status}:     & Verified against RTL testbench vectors \\
\textbf{Source}:     & \texttt{sw/reference\_model/} \\
\end{tabular}\\[1em]
\rule{0.7\textwidth}{0.4pt}
\vfill
\end{center}
\clearpage

%==============================================================================
\tableofcontents
\clearpage

%==============================================================================
\section{Overview}
\label{sec:overview}

Bit-accurate reference implementations of the KV cache engine in three
languages, all verified against the same test vectors.

\begin{table}[h]
\centering
\small
\begin{tabular}{@{}l l l l l@{}}
\toprule
Block & C++ class & \texttt{extern "C"} API & Python class & Tests \\
\midrule
\textbf{KV Cache Engine} & \texttt{lhsi::KVCacheEngine} & \texttt{lhsi\_kv\_*} & \texttt{KVCacheEngine} & 120 Python + 64 C++ \\
\bottomrule
\end{tabular}
\end{table}

\noindent Top-level orientation for the compiler team: see
\path{sw/README.md}. ISA specification: see
\path{docs/isa/kv_cache_engine_isa.pdf}.

%==============================================================================
\section{Building and Testing}
\label{sec:build}

\begin{lstlisting}[language=bash]
make test-all      # build C++ tests, run them, run all Python tests
make test          # only the C++ tests
make test-py       # only the Python tests
make shared        # libkv_cache_engine_ref.so
make static        # libkv_cache_engine_ref.a
make clean
\end{lstlisting}

\noindent Requires Python 3.10+, NumPy, and a C++17 compiler (gcc-9+ or
clang-10+). The C++ tests pick up the RTL test vectors from
\path{rtl/tb/testvectors/*.hex}; the Makefile regenerates them from
\path{analysis/gen_kv_testvectors.py} on first build if absent.

%==============================================================================
\section{KV Cache Engine}
\label{sec:kvcache}

\subsection{Configuration}

\begin{lstlisting}[language=C++]
#include "kv_cache_engine_ref.hpp"

lhsi::KVCacheEngineInfo info;
info.vector_dim    = 64;    // power of two (WHT constraint)
info.pq_bits       = 3;     // Lloyd-Max quantization bits
info.qjl_bits      = 1;     // QJL projection bits
info.sram_depth    = 1024;
info.coord_width   = 16;    // fixed-point coordinate width
info.coord_frac    = 12;    // fractional bits
info.norm_width    = 16;    // norm storage width
info.rotation_seed = 42;    // deterministic LCG seed
info.validate();            // asserts on invalid combinations
\end{lstlisting}

\subsection{C++ Class API}

\begin{lstlisting}[language=C++]
#include "kv_cache_engine_ref.hpp"

lhsi::KVCacheEngine engine;   // default config

// Compress a key vector
std::vector<int16_t> key_vec(64, /* ... */);
lhsi::CompressedKey ck = engine.compress_key(key_vec);

// Decompress back
std::vector<int16_t> reconstructed = engine.decompress_key(ck);

// Compress a value vector (no QJL)
lhsi::CompressedValue cv = engine.compress_value(val_vec);
std::vector<int16_t> val_recon = engine.decompress_value(cv);

// Store/load with SRAM address
engine.store_kv(/*addr=*/0, key_vec.data(), val_vec.data(), 64);
engine.load_k(0, key_out.data(), 64);
engine.load_v(0, val_out.data(), 64);
\end{lstlisting}

\subsection{Sub-Operations (Pipeline Stage Access)}

For compiler backends that need individual pipeline stages:

\begin{lstlisting}[language=C++]
// Each method mirrors one RTL pipeline stage
uint32_t norm = engine.compute_norm(vec, dim);
engine.normalize(vec, norm, normalized, dim);
engine.rotate(normalized, rotated, dim);
engine.quantize_pq(rotated, indices, dim);
engine.dequantize_pq(indices, dequantized, dim);
engine.qjl_compress(residual, signs, &res_norm, dim);
engine.qjl_decompress(signs, res_norm, correction, dim);
engine.inverse_rotate(corrected, output, dim);
\end{lstlisting}

\subsection{Plain C (\texttt{extern "C"} API)}

\begin{lstlisting}[language=C]
#include "kv_cache_engine_ref.hpp"

// Create with default config
lhsi_kv_handle_t* h = lhsi_kv_create_default();

// Or with custom config
lhsi_kv_info_t info = { .vector_dim = 64, .pq_bits = 3, /* ... */ };
lhsi_kv_handle_t* h = lhsi_kv_create(&info);

// Query config
lhsi_kv_info_t out_info = lhsi_kv_info(h);

// Compress key
uint32_t norm, res_norm;
uint8_t indices[64], signs[64];
lhsi_kv_compress_key(h, vec, 64, &norm, indices, &res_norm, signs);

// Decompress key
int16_t output[64];
lhsi_kv_decompress_key(h, norm, indices, res_norm, signs, 64, output);

// Compress/decompress value (no QJL fields)
lhsi_kv_compress_value(h, vec, 64, &norm, indices);
lhsi_kv_decompress_value(h, norm, indices, 64, output);

// Store/load with SRAM
lhsi_kv_store(h, 0, key_vec, val_vec, 64);
lhsi_kv_load_k(h, 0, key_out, 64);
lhsi_kv_load_v(h, 0, val_out, 64);

lhsi_kv_destroy(h);
\end{lstlisting}

\subsection{Python}

\begin{lstlisting}[language=Python]
from kv_cache_engine_ref import KVCacheEngine, KVCacheEngineInfo

engine = KVCacheEngine()

# Compress key
ck = engine.compress_key(key_vector)
reconstructed = engine.decompress_key(ck)

# Compress value
cv = engine.compress_value(val_vector)
val_recon = engine.decompress_value(cv)

# Store/load with SRAM
engine.store_kv(addr=0, key=key_vector, value=val_vector)
k_out = engine.load_k(addr=0)
v_out = engine.load_v(addr=0)

# Static method (no state needed)
ck = KVCacheEngine.decide_compress_key(key_vector)
\end{lstlisting}

%==============================================================================
\section{Numerical Semantics (Frozen)}
\label{sec:semantics}

All arithmetic is \textbf{fixed-point integer}, matching the RTL exactly:

\begin{itemize}
\item \textbf{Coordinates}: \texttt{COORD\_WIDTH}-bit signed two's complement,
      \texttt{COORD\_FRAC} fractional bits (default: Q4.12).
\item \textbf{Norms}: \texttt{NORM\_WIDTH}-bit unsigned, \texttt{NORM\_FRAC}
      fractional bits (default: Q8.8).
\item \textbf{WHT butterfly}: additions widen by 1 bit per stage
      ($\log_2 D$ stages), then truncate back to coordinate width.
\item \textbf{Quantization}: nearest-centroid with 3 comparators per
      coordinate. Centroids are Lloyd-Max optimal for $\mathcal{N}(0, 1/\sqrt{D})$.
\item \textbf{QJL}: random $\pm 1$ matrix from deterministic LCG seed.
      Inner product → sign extraction. Decompression scales by
      $\sqrt{\pi/2}$ (fixed-point constant).
\item \textbf{Integer sqrt}: non-restoring algorithm, bit-serial from MSB.
\end{itemize}

\noindent If the model disagrees with the chip on any input, the model is
wrong. Open an issue with the failing vector.

%==============================================================================
\section{Compressed Data Structures}
\label{sec:structures}

\subsection{\texttt{CompressedKey}}

\begin{table}[h]
\centering
\small
\begin{tabular}{@{}l l l p{6cm}@{}}
\toprule
Field & Type & Bits & Purpose \\
\midrule
\texttt{norm} & \texttt{uint32\_t} & 16 & L2 norm of original vector \\
\texttt{indices} & \texttt{uint8\_t[D]} & $3D$ & Lloyd-Max centroid indices (0--7) \\
\texttt{residual\_norm} & \texttt{uint32\_t} & 16 & L2 norm of quantization residual \\
\texttt{signs} & \texttt{uint8\_t[D]} & $D$ & QJL sign bits \\
\midrule
\multicolumn{2}{@{}l}{\textbf{Total (D=64)}} & \textbf{288} & 3.56$\times$ compression \\
\bottomrule
\end{tabular}
\end{table}

\subsection{\texttt{CompressedValue}}

\begin{table}[h]
\centering
\small
\begin{tabular}{@{}l l l p{6cm}@{}}
\toprule
Field & Type & Bits & Purpose \\
\midrule
\texttt{norm} & \texttt{uint32\_t} & 16 & L2 norm of original vector \\
\texttt{indices} & \texttt{uint8\_t[D]} & $3D$ & Lloyd-Max centroid indices (0--7) \\
\midrule
\multicolumn{2}{@{}l}{\textbf{Total (D=64)}} & \textbf{208} & 4.92$\times$ compression \\
\bottomrule
\end{tabular}
\end{table}

%==============================================================================
\section{Verification Status}
\label{sec:verification}

\begin{table}[h]
\centering
\small
\begin{tabular}{@{}l l l@{}}
\toprule
Test & Block / Aspect & Status \\
\midrule
120 Python unit tests & Full round-trip, all paths & Pass \\
64 C++ unit tests & Full round-trip, all paths & Pass \\
Compress/decompress round-trip & K and V paths & Bit-exact \\
K/V asymmetry verification & QJL on K only & Pass \\
Streaming vs batch agreement & tick() vs compress\_*() & Pass \\
\texttt{extern "C"} = C++ class & C API parity & Pass \\
RTL replay vectors & Hex-file bit-exact match & Pass \\
Compression ratio >= 3$\times$ & Canonical trace & Pass \\
\bottomrule
\end{tabular}
\end{table}

%==============================================================================
\section{Co-Design Context}
\label{sec:codesign}

This directory is the deliverable for compiler-team integration.
The reference model provides the same API contract that the hardware
will expose, enabling compiler backends to target the KV cache engine
before silicon or FPGA prototypes exist.

\medskip
\noindent Next integration steps:
\begin{enumerate}
\item Compiler targets Python/C++ reference models (Phase~0 --- now)
\item AXI-Lite + AXI-Stream on Zynq UltraScale+ FPGA (Phase~1)
\item Multi-block FPGA project with precision controller (Phase~2)
\item TSMC 16FFC silicon (Phase~3, 2027+)
\end{enumerate}

\vspace{0.5in}
\noindent\rule{\textwidth}{0.4pt}\\
\noindent\small\textit{LonghornSilicon KV Cache Engine --- Reference Model API.
This document is open and may be redistributed.}

\end{document}
